% problem-000210 3 硼硅低配位化学
\section*{3. 硼硅低配位化学}
甲酸受热分解有脱⽔和脱羧两种⽅式：⽅式 (1) $\mathrm { H C O O H } = \mathrm { H _ { 2 } O } + \mathrm { C O }$ （脱⽔）⽅式 (2) $\mathrm { H C O O H } = \mathrm { H } _ { 2 } + \mathrm { C O } _ { 2 }$ （脱羧）

在⾼温⽔热⽆催化剂条件下，主要发⽣⽅式(1)。为获得⽔热条件下脱⽔反应的相关数据，开展了以下实验：室温下向耐压⽯英核磁管中加⼊1mol L−1的甲酸⽔溶液并密封，在温度�下反应⼀段时间后，记录温度 � 下核磁管中液相的体积占⽐ $x _ { T }$ ，随后将⽯英管快速降⾄室温终⽌反应（此时液相占⽐与反应温度�时的不同），利⽤核磁分析液相和⽓相中各物种的浓度，通过计算得到甲酸分解的转化率，进⽽分析反应热⼒学和动⼒学。甲酸⽔热脱羧反应的机理如图3.1所⽰，其中质⼦化甲酸中间体的⽣成是快速平衡反应：

（图：）
图 3.1 甲酸⽔热脱⽔反应的机理

研究表明，当向反应体系中加⼊HCl，反应速率 $r _ { 1 }$ 与甲酸和H+ 浓度的关系为：

$
r _ {1} = k _ {1} (\mathrm{HCOOH}) \left(\mathrm{H} ^ {+}\right)
$

若反应体系中不加HCl，甲酸分解速率 $r _ { 2 }$ 与甲酸浓度的关系为：

$
r _ {2} = k _ {2} (\mathrm{HCOOH}) ^ {1. 5}
$

\subsection*{3-1}
推导加HCl和不加HCl两种情形下的反应动⼒学⽅程。

\subsection*{3-2}
通过反应速率常数 $k _ { 1 }$ 和 $k _ { 2 }$ 可以获得甲酸在⽔热条件下的酸解平衡常 $K _ { a }$ ，在 $2 0 0 { \sim } 3 0 0 ^ { \circ } \mathrm { C }$ 范围内，测得不同温度下的反应速率常数 $k _ { 1 }$ 和 $k _ { 2 }$ ⻅下表：

\textit{（原卷此处含表格/仪器试剂表，详见 PDF 或 MD 源文件。）}

\subsection*{3-2}
-1 推导 $\mathsf { p } K _ { a }$ 与 $k _ { 1 }$ 和 $k _ { 2 }$ 的关系式。

\subsection*{3-2}
-2 计算上表中不同温度下的 $\mathsf { p } K _ { a }$ ，并解释其随温度变化的原因。

\subsection*{3-3}
甲酸⽔热脱羧反应的平衡常数�定义如下：

$
K = [ \mathrm{CO} (\mathrm{aq}) ] [ \mathrm{HCOOH} (\mathrm{aq}) ]
$

其中，[CO(aq)]和[HCOOH(aq)]分别为达平衡时液相中CO和甲酸的浓度。体系中CO在⽓相 $( \mathrm { g } )$ 和液相 (aq)中都存在，⽽甲酸主要存在于液相中。本实验需将反应体系冷⾄室温后才能进⾏产物浓度分析，但⽆法直接测得反应温度�下的液相浓度，测得的是整个反应体系中（包含⽓相和液相）达到平衡后 CO 和甲酸的摩尔数之⽐：$Q \ = \ n ( \mathrm { C O } ) / n ( \mathrm { H C O O H } )$ 。CO在液相和⽓相中的分配系数 $K _ { D }$ 定义为： $K _ { D } = [ \mathrm { C O ( g ) / C O ( a q ) } ]$

推导�的表达式，⽤�、 $K _ { D }$ 和 $x _ { T }$ 等实验中的可观测量来表达。

\subsection*{3-4}
实验数据表明，在 $2 0 0 { \sim } 3 0 0 ^ { \circ } \mathrm { C }$ 的温度范围内，lg�和lg $K _ { D }$ 均与温度的倒数 $1 / T$ 呈较好的线性关系，如下式所⽰，其中�以K为单位：

$
\lg K = - 3. 0 \times 1 0 ^ {3} / T + 5. 9; \quad \lg K _ {D} = 1. 3 \times 1 0 ^ {3} / T - 1. 4
$

计算 $2 0 0 { \sim } 3 0 0 \ { } ^ { \circ } \mathrm { C }$ 内甲酸脱⽔反应的标准摩尔焓变 $\Delta H ^ { \ominus } \mathrm { _ { C O } }$ 和CO在液相中的溶解焓 $\cdot _ { \Delta H } \Theta _ { \mathrm { p } }$ C

\subsection*{3-5}
实验发现，当 $x _ { T } < 0 . 1$ 时，在 $2 0 0 { \sim } 3 0 0 \ { } ^ { \circ } \mathrm { C }$ 范围内，lg�与 $1 / T$ 也具有良好的线性关系，其斜率为 $k _ { Q }$ 满⾜以下关系：−2.303 � $k _ { Q } = \Delta H ^ { \ominus } \mathrm { c o } + \beta \Delta H ^ { \ominus } \mathrm { _ { D c } }$

\subsection*{3-5}
-1 推导 $\beta$ 的表达式，⽤ $K _ { D }$ 和 $x _ { T }$ 表⽰。

\subsection*{3-5}
-2 解释当 $x _ { T } < 0 . 1$ 时，lg � 与 $1 / T$ 具有良好线性关系的原因。

提⽰：对于函数 $y = f ( x )$ ，� 对 � 的导数 d $y /  { \mathrm { d } } x$ 如下表所⽰（表中�和�均为常数，�是�的函数）：

\textit{（原卷此处含表格/仪器试剂表，详见 PDF 或 MD 源文件。）}

\subsection*{3-6}
选择合适的催化剂，可使甲酸在温和条件下⾼选择性发⽣脱羧反应，获得氢⽓。Ru配合物Ru-1就是其中之⼀，其催化甲酸脱羧的机理如图3.2所⽰（注：循环箭头上只给出了部分反应物和离去物）。

（图：）
图 3.2 过渡⾦属催化的甲酸脱羧反应

图 3.2 中，A、B、C 均为 Ru 的 6 配位化合物，其中均含有 Cl 元素；A′和 C′的结构类似于 A 和 $\mathbf { C } ,$ ，但不含氯；X为某种反应物；Ru的价态在催化循环中保持不变。写出配合物A、B、C、A′、C′的结构（有机配体可以⽤ $\mathrm { N P _ { 3 } }$ 代表）。
